\documentclass[11pt,a4paper]{report}
\usepackage[margin=1in]{geometry}
\usepackage{amsmath,amssymb,mathtools,bm}
\usepackage{booktabs}
\usepackage{array}
\usepackage{longtable}
\usepackage{enumitem}
\usepackage{hyperref}
\usepackage{xcolor}
\usepackage{listings}
\usepackage{microtype}
\usepackage[T1]{fontenc}
\usepackage[utf8]{inputenc}

\hypersetup{colorlinks=true,linkcolor=blue,citecolor=blue,urlcolor=blue}
\lstset{basicstyle=\ttfamily\small,breaklines=true,columns=fullflexible,frame=single}

\newcommand{\dd}{\mathrm{d}}
\newcommand{\ii}{\mathrm{i}}
\newcommand{\eps}{\varepsilon}
\newcommand{\LTD}{\operatorname{LTD}}
\newcommand{\CFF}{\operatorname{CFF}}
\newcommand{\OSE}{\operatorname{OSE}}
\newcommand{\Res}{\operatorname{Res}}
\newcommand{\Graph}{\mathcal G}
\newcommand{\Basis}{\mathcal B}
\newcommand{\Orient}{\mathcal O}
\newcommand{\Surf}{\mathcal S}
\newcommand{\Weights}{\mathcal W}
\newcommand{\Active}{\mathcal A}
\newcommand{\Samples}{\mathcal X}
\newcommand{\Ints}{\mathcal E}
\newcommand{\Rep}{\mathcal R}
\newcommand{\Pinch}{\mathcal P}
\newcommand{\rank}{\operatorname{rank}}

\title{Hybrid Three-Dimensional LTD/CFF Representations}
\author{Implementation and theory note}
\date{April 2026}

\begin{document}
\maketitle

\begin{abstract}
This note describes the derivative-free three-dimensional representation
implemented in \texttt{hybrid3d.py}.  The construction combines loop-tree
duality (LTD) with the causal/CFF representation in order to handle repeated
propagators while keeping the numerator as a sampled black box.  It also
describes the bounded-energy-degree extension used when edge-momentum
representation (EMR) numerators contain powers above the affine CFF class.

The document is organized as follows.  Chapter~1 gives the mathematical
construction, with all dependencies made explicit.  Chapter~2 explains how the
Python implementation and the oriented JSON schema encode the formulae.
Chapter~3 reports current Symbolica evaluator timings from the example scripts.
\end{abstract}

\tableofcontents

\chapter{Theoretical Construction}

\section{Framing, Notation, LTD, and CFF}

Let \(\Graph\) be a connected scalar graph with \(L\) loop momenta and internal
edge set \(E\).  A loop-momentum basis (LMB) is denoted by
\(\Basis=(k_1,\ldots,k_L)\).  For an edge \(e\in E\),
\begin{equation}
  q_e(k,p)=\sum_{r=1}^{L}\rho_{er}k_r+\sum_{a=1}^{n_{\rm ext}}\eta_{ea}p_a,
  \qquad \rho_{er},\eta_{ea}\in\mathbb Z ,
\end{equation}
and
\begin{equation}
  q_e^0(k^0,p^0)=\rho_e\cdot k^0+\eta_e\cdot p^0,\qquad
  \omega_e(\vec k,p,m)=\sqrt{\vec q_e^{\,2}+m_e^2}.
\end{equation}
The Feynman denominator is
\begin{equation}
  D_e=(q_e^0)^2-\omega_e^2+\ii0 .
\end{equation}

For a numerator \(N\) in EMR variables, the \(D\)-dimensional input integral is
\begin{equation}
  I_{\Graph,\Basis}[N]
  =
  \int_{\mathbb R^{4L}}
  \prod_{r=1}^{L}\frac{\dd^4 k_r}{(2\pi)^4}\,
  N(\{q_e(k,p)\}_{e\in E})
  \prod_{e\in E}D_e^{-a_e}.
  \label{eq:four-dimensional-integral}
\end{equation}
The powers \(a_e\) are normally one; repeated propagators are encoded below as
explicit repeated edges or, equivalently, as collapsed channels with powers.

\paragraph{LTD dependency data.}
The LTD construction is an algorithmic residue expansion of
\eqref{eq:four-dimensional-integral} after choosing
\begin{itemize}[leftmargin=2em]
  \item the graph \(\Graph\),
  \item the loop-momentum basis \(\Basis\),
  \item a contour-closure prescription \(\kappa\), i.e. the side on which each
        loop-energy integration is closed.
\end{itemize}
We denote the resulting spatial integrand by
\begin{equation}
  I^{(3)}_{\LTD}[\Graph,\Basis,\kappa;N](\vec k;p,m)
  =
  \sum_{\mathfrak o\in\Orient_{\LTD}(\Graph,\Basis,\kappa)}
  N\!\left(Q_{\mathfrak o}^{\LTD}(\vec k;p,m)\right)
  W_{\mathfrak o}^{\LTD}(\vec k;p,m).
\end{equation}
The map \(Q_{\mathfrak o}^{\LTD}\) gives every EMR edge energy after solving
the cut equations.  The weight \(W_{\mathfrak o}^{\LTD}\) is a product of
cut half-edge factors and LTD dual surfaces.  In this note we do not rederive
LTD; we use the multiloop algorithmic construction of
Ref.~\cite{ltd1906}.

\paragraph{CFF dependency data.}
For CFF/causal representations we use the graph as the primary input and assume
the canonical source/sink contraction ordering used by the implementation.
Thus
\begin{equation}
  I^{(3)}_{\CFF}[\Graph;N](\vec k;p,m)
  =
  \sum_{\mathfrak o\in\Orient_{\CFF}(\Graph)}
  N\!\left(Q_{\mathfrak o}^{\CFF}(\vec k;p,m)\right)
  W_{\mathfrak o}^{\CFF}(\vec k;p,m).
\end{equation}
The CFF weights use causal \(E\)-surfaces in the denominator.  The construction
is the one of Refs.~\cite{cff2211,cff2311}.  The standard black-box numerator
assumption is that \(N\) is at most affine in each EMR edge energy \(q_e^0\).

\paragraph{Surface notation.}
For any integer vector \(c_e\in\{-1,0,1\}\) and external-energy shift
\(\Delta\), define the signed causal surface
\begin{equation}
  \eta(c,\Delta)=\sum_{e}c_e\,\omega_e+\Delta .
\end{equation}
An \(E\)-surface has all non-zero internal-energy coefficients with the same
orientation after the canonical CFF construction.  An \(H\)-surface is a more
general signed hyperboloid surface.  Pure CFF denominator factors use only
\(E\)-surfaces; LTD and hybrid terms may use \(H\)-surfaces when the residue
geometry requires them.  Numerator-side surfaces are allowed in the numerator
and do not introduce denominator singularities.

\section{Repeated Propagators Without Numerator Derivatives}

Repeated propagators correspond to groups
\begin{equation}
  R_g=\{r_{g,1},\ldots,r_{g,\nu_g}\},\qquad \nu_g\ge2,
\end{equation}
whose members have the same canonical routing and mass, possibly with reversed
edge orientation.  Let \(Q_g^0\) be the canonical repeated-channel energy and
\(\omega_g\) the common on-shell energy.  The purpose of the hybrid
construction is to evaluate the repeated-channel residues without applying
derivatives to the black-box numerator.

\phantomsection
\subsection*{1.2.a The Fully Generic Formula}
\addcontentsline{toc}{subsection}{1.2.a The Fully Generic Formula}

Start from
\begin{align}
  I_{\Graph,\Basis}[N]
  &=
  \int \prod_{r=1}^{L}\frac{\dd^4 k_r}{(2\pi)^4}\,
  N(\{q_e(k,p)\})
  \prod_{s\in S}D_s^{-1}
  \prod_{g\in G_{\rm rep}}\prod_{a=1}^{\nu_g}D_{r_{g,a}}^{-1}
  \notag\\
  &=
  \int \prod_{r=1}^{L}\frac{\dd^4 k_r}{(2\pi)^4}\,
  N(\{q_e(k,p)\})
  \prod_{s\in S}D_s^{-1}
  \prod_{g\in G_{\rm rep}}D_g^{-\nu_g}.
  \label{eq:collapsed-repeated-integral}
\end{align}
The second equality is the collapsed denominator identity.  It is only a
denominator identity: the numerator still has one EMR argument for every
original edge.

Let \(\bar E=S\sqcup G_{\rm rep}\) be the collapsed edge set and assign powers
\[
  n_{\bar e}=
  \begin{cases}
    1, & \bar e=s\in S,\\
    \nu_g, & \bar e=g\in G_{\rm rep}.
  \end{cases}
\]
For every rank-\(L\) collapsed cut basis \(B\subset\bar E\) and pole-sign vector
\(\sigma_B\), solve the cut equations
\begin{equation}
  Q_b^0(k^0,p^0)=\sigma_b\omega_b,\qquad b\in B,
  \label{eq:cut-equations}
\end{equation}
for the loop energies:
\begin{equation}
  k^0=K_{B,\sigma_B}(\vec k;p,m).
\end{equation}
The direct residue theorem gives the strict equality
\begin{align}
  I_{\Graph,\Basis}[N]
  &=
  \int \prod_{r=1}^{L}\frac{\dd^3 k_r}{(2\pi)^3}
  \sum_{B,\sigma_B}
  C_{B,\sigma_B}
  \Biggl[
  \prod_{b\in B}
  \frac{1}{(n_b-1)!}
  \left(\frac{\partial}{\partial Q_b^0}\right)^{n_b-1}
  \notag\\[-2mm]
  &\hspace{34mm}
  \frac{
    N(\{q_e^0(k^0,p^0),\vec q_e\})
  }{
    \prod_{b\in B}(Q_b^0+\sigma_b\omega_b)^{n_b}
    \prod_{\bar e\notin B}D_{\bar e}^{n_{\bar e}}
  }
  \Biggr]_{k^0=K_{B,\sigma_B}} .
  \label{eq:raw-higher-residue}
\end{align}
The constant \(C_{B,\sigma_B}\) contains the orientation sign and the Jacobian
of the chosen energy basis.

Equation~\eqref{eq:raw-higher-residue} is not yet acceptable when a powered
channel is cut, because the derivatives act on \(N\).  Define the set of active
original edges
\begin{equation}
  \Active_B
  =
  \left\{
  e\in E\;:\;
  q_e^0 \hbox{ depends on } Q_b^0
  \hbox{ for some } b\in B \hbox{ with } n_b>1
  \right\}.
  \label{eq:active-edges}
\end{equation}
For the affine EMR class, insert the exact interpolation identity
\begin{equation}
  N(\{q_e^0\})
  =
  \sum_{\theta\in\{\pm1\}^{\Active_B}}
  N\!\left(\{q_{e,\theta}^0\}\right)
  \prod_{e\in\Active_B}
  \frac{1+\theta_e q_e^0/\omega_e}{2},
  \label{eq:affine-interpolation}
\end{equation}
with \(q_{e,\theta}^0=\theta_e\omega_e\) for active edges and the ordinary
affine value for inactive edges.  In the physical CFF-localized form, the same
interpolation is organized as a finite sum over repeated-copy sign orbits
obtained by solving the localizing delta constraints before taking the
equal-energy limit.  Both forms are exact for the affine EMR numerator class.

Substituting \eqref{eq:affine-interpolation} into
\eqref{eq:raw-higher-residue} gives
\begin{align}
  I_{\Graph,\Basis}[N]
  &=
  \int \prod_{r=1}^{L}\frac{\dd^3 k_r}{(2\pi)^3}\,
  I^{\rm HybridLTD}_{\Graph,\Basis,\kappa}[N](\vec k;p,m),
  \label{eq:spatial-hybrid-integral}\\
  I^{\rm HybridLTD}_{\Graph,\Basis,\kappa}[N]
  &=
  \sum_{\mathfrak o\in\Orient_{\rm H}(\Graph,\Basis,\kappa)}
  N\!\left(Q_{\mathfrak o}^{\rm H}(\vec k;p,m)\right)
  W_{\mathfrak o}^{\rm H}(\vec k;p,m).
  \label{eq:hybrid-master}
\end{align}
Here one orientation is
\[
  \mathfrak o=(B,\sigma_B,\theta,\mu),
\]
where \(\mu\) labels one algebraic term from expanding the known derivative in
\eqref{eq:raw-higher-residue}.  The numerator is never differentiated; it is
only sampled at \(Q_{\mathfrak o}^{\rm H}\).  The weight is a finite product
\begin{equation}
  W_{\mathfrak o}^{\rm H}
  =
  c_{\mathfrak o}
  \prod_{h\in H_{\mathfrak o}}\frac{1}{2\omega_h}
  \prod_{\eta\in D_{\mathfrak o}}\eta(\vec k;p,m)^{-m_{\mathfrak o,\eta}}
  \prod_{\zeta\in N_{\mathfrak o}}\zeta(\vec k;p,m) .
  \label{eq:hybrid-weight}
\end{equation}
The denominator factors \(D_{\mathfrak o}\) are cached \(E\)- or \(H\)-surfaces;
the numerator factors \(N_{\mathfrak o}\) are known surface factors produced by
interpolation and contact reductions.  If no repeated channel is present, then
\(n_b=1\) for all cuts, \(\Active_B=\varnothing\), and
\eqref{eq:hybrid-master} collapses exactly to the LTD representation with the
same \((\Graph,\Basis,\kappa)\).

\phantomsection
\subsection*{1.2.b The Box with a Triple Repeated Propagator}
\addcontentsline{toc}{subsection}{1.2.b The Box with a Triple Repeated Propagator}

Consider the one-loop box routing
\[
  q_0=k,\qquad q_1=k+p_1,\qquad q_2=k+p_1+p_2,\qquad
  Q=k+p_1+p_2+p_3,
\]
with three repeated copies of \(Q\).  The collapsed denominator is
\[
  D_0D_1D_2D_Q^3 .
\]
Write \(\Delta_i=p_1^0+\cdots+p_i^0\) and define the causal surfaces
\begin{align}
  E_{ij}^{\sigma_i\sigma_j}
  &=
  \sigma_i\omega_i+\sigma_j\omega_j+\Delta_i-\Delta_j,\\
  E_{iQ}^{\sigma_i\sigma_Q}
  &=
  \sigma_i\omega_i+\sigma_Q\omega_Q+\Delta_i-\Delta_3 .
\end{align}
The three simple-cut sectors are
\begin{equation}
  I_i^{\rm simple}
  =
  -\frac{
    N\!\left(
      q_0^0,q_1^0,q_2^0,Q^0,Q^0,Q^0
    \right)_{k^0=-\Delta_i+\omega_i}
  }{
    2\omega_i\,
    E_{ij}^{++}\,
    E_{i\ell}^{++}\,
    \left(E_{iQ}^{++}\right)^3
  },
  \qquad \{i,j,\ell\}=\{0,1,2\}.
  \label{eq:box-simple-cuts}
\end{equation}
The signs shown correspond to the upper on-shell pole convention used in the
examples; the opposite closure is obtained by replacing every displayed
\(+\omega\) by the corresponding signed on-shell energy.

For the repeated cut, set \(Q^0=\sigma\omega_Q\) and
\[
  k^0=\sigma\omega_Q-\Delta_3 .
\]
Let \(\theta=(\theta_3,\theta_4,\theta_5)\) be the sample signs for the three
repeated EMR copies.  The numerator call is
\begin{equation}
  N_{\sigma,\theta}
  =
  N\!\left(
    \sigma\omega_Q-\Delta_3,\;
    \sigma\omega_Q-\Delta_2,\;
    \sigma\omega_Q-\Delta_1,\;
    \theta_3\omega_Q,\;
    \theta_4\omega_Q,\;
    \theta_5\omega_Q
  \right).
  \label{eq:box-repeated-samples}
\end{equation}
The repeated-cut contribution has the finite form
\begin{equation}
  I_Q^{\rm rep}
  =
  \sum_{\sigma=\pm1}
  \sum_{\theta\in\{\pm1\}^3}
  N_{\sigma,\theta}
  \sum_{\mu}
  c_{\sigma,\theta,\mu}
  \frac{
    \prod_{\zeta\in N_{\sigma,\theta,\mu}}\zeta
  }{
    (2\omega_Q)^{3+r_\mu}
    \prod_{i=0}^{2}
    \left(E_{Qi}^{\sigma+}\right)^{m_{\mu i}}
  } .
  \label{eq:box-repeated-explicit}
\end{equation}
All objects in \eqref{eq:box-repeated-explicit} are explicit functions of
external energies and on-shell energies.  The coefficients
\(c_{\sigma,\theta,\mu}\), powers \(r_\mu,m_{\mu i}\), and numerator-side
surface list \(N_{\sigma,\theta,\mu}\) are the expansion of the second-order
confluent residue.  In the JSON they appear as rational variant prefactors,
repeated half-edge entries for \(2\omega_Q\), and repeated surface ids.  The
final boxed formula is
\[
  I^{\rm HybridLTD}_{\Box^{(3)}} =
  I_0^{\rm simple}+I_1^{\rm simple}+I_2^{\rm simple}+I_Q^{\rm rep}.
\]

\section{Higher Energy Powers in Numerators}

The CFF construction is naturally compatible with numerators affine in each
EMR edge energy.  When a user supplies a bound \(d_e\) for the highest power of
each edge energy in \(N\), the implementation first checks loop-energy UV
convergence.  In every loop-energy direction \(v\), the degree of the numerator
minus twice the number of denominators depending on that direction must be
strictly smaller than \(-1\).  If this check fails, finite pole data are not
enough and the command refuses to build a representation.

\phantomsection
\subsection*{1.3.1.a Quadratic Bounds}
\addcontentsline{toc}{subsection}{1.3.1.a Quadratic Bounds}

Let \(x=q_e^0\), \(\omega=\omega_e\), and suppose \(N\) has degree at most two
in \(x\).  There is an exact decomposition
\begin{equation}
  N(x)=R_eN(x)+\left(x^2-\omega^2\right)C_eN,
  \label{eq:quadratic-division}
\end{equation}
where the pole-safe affine remainder is
\begin{equation}
  R_eN(x)=
  \frac{x+\omega}{2\omega}N(+\omega)
  +
  \frac{\omega-x}{2\omega}N(-\omega),
  \label{eq:quadratic-remainder}
\end{equation}
and the contact quotient is
\begin{equation}
  C_eN=
  \frac{N(+\omega)-2N(0)+N(-\omega)}{2\omega^2}.
  \label{eq:quadratic-contact}
\end{equation}
Therefore
\begin{equation}
  \frac{N}{D_e\prod_{f\ne e}D_f}
  =
  \frac{R_eN}{D_e\prod_{f\ne e}D_f}
  +
  \frac{C_eN}{\prod_{f\ne e}D_f}.
  \label{eq:quadratic-master}
\end{equation}
The first term is affine in \(q_e^0\) and can be handled by ordinary CFF.  The
second term is a pinched graph with edge \(e\) removed, also handled by CFF.
Consequently the pure CFF denominator remains \(E\)-surface only.  Known
factors such as \(x+\omega\), \(\omega-x\), and \(1/\omega^2\) are stored as
numerator-side surface factors or half-edge powers; they are not denominator
surfaces.  In hybrid mode the same \(R/C\) reduction is applied before the
remaining repeated-channel confluent residue is expanded, so the black-box
numerator is still sampled only at explicit energies.

\phantomsection
\subsection*{1.3.1.b One-Loop Box with One Quadratic Edge}
\addcontentsline{toc}{subsection}{1.3.1.b One-Loop Box with One Quadratic Edge}

For a non-repeated one-loop box with denominators \(D_0D_1D_2D_3\), assume the
only non-affine dependence is \(d_0=2\).  Equations
\eqref{eq:quadratic-division}--\eqref{eq:quadratic-master} give
\begin{equation}
  I_{\Box}^{\CFF}[N]
  =
  I_{\CFF}[\Box;R_0N]
  +
  I_{\CFF}[\Box/0;C_0N],
  \label{eq:box-quadratic}
\end{equation}
where \(\Box/0\) is the triangle obtained by deleting propagator \(0\).  More
explicitly,
\begin{align}
  I_{\CFF}[\Box;R_0N]
  &=
  \sum_{\mathfrak o\in\Orient_{\CFF}(\Box)}
  \left[
  \frac{q_0^0(\mathfrak o)+\omega_0}{2\omega_0}N(q_0^0=+\omega_0)
  +
  \frac{\omega_0-q_0^0(\mathfrak o)}{2\omega_0}N(q_0^0=-\omega_0)
  \right]
  W_{\mathfrak o}^{\CFF},\\
  I_{\CFF}[\Box/0;C_0N]
  &=
  \sum_{\mathfrak t\in\Orient_{\CFF}(\Box/0)}
  \frac{
    N(+\omega_0)-2N(0)+N(-\omega_0)
  }{2\omega_0^2}
  W_{\mathfrak t}^{\CFF}.
\end{align}
Both \(W_{\mathfrak o}^{\CFF}\) and \(W_{\mathfrak t}^{\CFF}\) contain only
\(E\)-surface denominators.

\phantomsection
\subsection*{1.3.2.a Bounds Above Quadratic}
\addcontentsline{toc}{subsection}{1.3.2.a Bounds Above Quadratic}

For \(d_e>2\), let \(a_1,\ldots,a_{d_e-1}\) be non-pole interpolation nodes
with \(a_j\ne\pm1\), and write \(x=q_e^0/\omega_e\).  If \(L_j(x)\) is the
Lagrange basis on these nodes, the contact quotient is reconstructed from
black-box samples by
\begin{equation}
  C_eN(q_e^0)
  =
  \sum_{j}
  \frac{L_j(q_e^0/\omega_e)}{(a_j^2-1)\omega_e^2}
  \left[
    N(a_j\omega_e)
    -\frac{a_j+1}{2}N(+\omega_e)
    -\frac{1-a_j}{2}N(-\omega_e)
  \right].
  \label{eq:higher-contact}
\end{equation}
The remainder part is still the pole interpolation at \(\pm\omega_e\).  The
new feature is that \(C_eN\) is now a known polynomial in the remaining energy
variables.  Those known factors must be carried as numerator-side surfaces and,
when their degree is still too high for CFF, recursively reduced by the same
finite-pole division.  This is why higher powers require \(E\)- and
\(H\)-surface factors in the numerator even though pure CFF denominator trees
remain \(E\)-surface only.  Hybrid mode uses the same bounded-degree contact
algebra before or inside each repeated-channel residue sector.

\phantomsection
\subsection*{1.3.2.b One-Loop Box with One Cubic Edge}
\addcontentsline{toc}{subsection}{1.3.2.b One-Loop Box with One Cubic Edge}

For \(N(x)=x^3 M\), with \(M\) independent of \(x=q_0^0\),
\[
  x^3=\omega_0^2 x+x(x^2-\omega_0^2).
\]
Thus
\begin{equation}
  \frac{x^3 M}{D_0D_1D_2D_3}
  =
  \frac{\omega_0^2 xM}{D_0D_1D_2D_3}
  +
  \frac{xM}{D_1D_2D_3}.
  \label{eq:box-cubic}
\end{equation}
The first term has an affine numerator on the box.  The second term is a
pinched triangle with a known numerator-side linear energy factor \(x\).  After
the triangle CFF construction samples its numerator map, \(x\) becomes a cached
surface factor in the numerator.  For multiloop examples the same mechanism can
produce numerator-side \(H\)-surfaces because a pinched lower-sector energy can
be a signed difference of on-shell energies.  These numerator factors are
regular polynomial factors, not denominator singularities.

\chapter{Python Implementation and JSON Format}

\section{Module Structure}

The implementation is organized around an evaluator-complete intermediate
format.
\begin{itemize}[leftmargin=2em]
  \item \texttt{graph\_io.py} parses DOT graphs, extracts momentum signatures,
        detects repeated channels, and builds split-mass reference graphs.
  \item \texttt{graph\_signatures.py} implements the
        \texttt{graph\_from\_signatures} command, which turns products such as
        \texttt{prop(k1+p1,m)} into DOT graphs.
  \item \texttt{orientation\_bundle.py} constructs LTD, CFF, hybrid, and
        bounded-degree contact bundles.  This is where repeated channels are
        collapsed, active numerator maps are enumerated, and CFF lower minors
        are generated.
  \item \texttt{structure.py} serializes bundles to schema-v6 JSON and contains
        the arbitrary-precision Python evaluator.
  \item \texttt{pretty.py} renders the JSON in a human-readable form.
  \item \texttt{symbolica\_eval.py} compiles a fixed-numerator JSON expression
        to Symbolica eager and assembly-backed evaluators.
\end{itemize}

\section{Oriented JSON Semantics}

The top-level JSON contains graph metadata, a surface cache, and a list of
orientations.  The central invariant is:
\begin{quote}
one orientation corresponds to one unique EMR edge-energy numerator map.
\end{quote}
All denominator terms sharing that numerator call are stored as
\texttt{variants}.  Each variant has its own rational prefactor, half-edge
list, numerator-surface list, and factorized denominator tree.

The evaluator uses only the JSON plus the DOT graph.  It does not reconstruct a
hidden LTD or CFF object.  A variant contributes
\[
  \texttt{pref}\;
  \prod_{e\in\texttt{half\_edges}}(2\omega_e)^{-1}
  \prod_{s\in\texttt{num\_surfaces}}\eta_s
  \prod_{t\in\texttt{denominator tree}}\eta_t^{-1}
\]
times the numerator sampled at the orientation's \texttt{edge\_q0} map.

\section{Surface Cache}

Each surface cache entry records a class, coefficients of internal on-shell
energies, and external-energy shifts.  The class is printed as
\texttt{e} or \texttt{h}.  If a surface appears only in numerator factors, the
entry carries \texttt{numerator\_only=true} and pretty printing displays
\texttt{(e)} or \texttt{(h)}.  If a numerator-side surface is identical to an
existing denominator surface, the same id is reused.

\section{Orientation Labels and Variants}

Compact labels use \(+\), \(-\), \(0\), and \(x\):
\begin{itemize}[leftmargin=2em]
  \item \(+\): the numerator samples that edge at \(+\omega_e\);
  \item \(-\): the numerator samples that edge at \(-\omega_e\);
  \item \(0\): the numerator samples that edge at zero energy;
  \item \(x\): the edge has a non-trivial affine map, printed by
        \texttt{--show-details}.
\end{itemize}
The physical origin of a contribution, such as CFF, LTD, pinched contacts, or
hybrid basis/interpolation data, is variant metadata rather than part of the
orientation identity.  This is what keeps the number of numerator calls
minimal.

\section{Command-Line Workflow}

The main commands are:
\begin{lstlisting}
python3 hybrid3d.py validate --dot examples/graphs/box_pow3.dot
python3 hybrid3d.py build --family hybrid --dot examples/graphs/box_pow3.dot --pretty
python3 hybrid3d.py test --dot examples/graphs/box_pow3.dot
python3 hybrid3d.py test-cff-ltd --dot examples/graphs/box.dot --energy-degree-bounds 0:2
python3 hybrid3d.py compile --orientation-json demo.json --dot graph.dot --numerator-expr "1"
python3 hybrid3d.py evaluate --orientation-json demo.json --dot graph.dot --evaluator-backend symbolica
\end{lstlisting}
The repeated-propagator three-way test compares pure CFF, hybrid, and the
split-mass LTD limiting proxy.  The ordinary LTD formula on an unsplit repeated
graph is not a valid physics comparison because numerator derivatives are not
included there.

\section{Symbolica Evaluators}

The \texttt{compile} command accepts one JSON representation and one fixed
numerator expression.  It reconstructs the JSON expression symbolically,
introducing structured functions such as \texttt{OSE(qx,qy,qz,m)},
\texttt{etaE(i)}, \texttt{etaH(i)}, and tagged edge/loop/external momentum
components.  The resulting Symbolica evaluator is saved both as a shared
library and as eager evaluator state.  The JSON is updated with relative paths
and fingerprints tying the evaluator to the graph, numerator, value type, and
orientation structure.

The \texttt{evaluate} command can use:
\[
\texttt{builtin},\quad
\texttt{symbolica\_compiled},\quad
\texttt{symbolica\_eager},\quad
\texttt{symbolica\_eager\_symjit}.
\]
Profiling performs ten calls of a batch evaluator.  The default batch size is
100, so a profile line reports one thousand repeated samples of the same input
kinematics.

\chapter{Performance}

\section{Benchmark Setup}

The numbers below were generated on April 26, 2026 using the scripts in
\texttt{examples/scripts}.  Each timing is the per-sample time reported by
\texttt{evaluate --profiling 100}, i.e. ten batch calls of size 100.  The
15-external one-loop benchmark is intentionally omitted here: the LTD build is
small, but the pure CFF one-loop representation has \(2^{15}-2=32766\) sign
sectors and the CFF build is not practical for routine documentation runs.

\section{Repeated Box}

The repeated box script
\texttt{examples/scripts/box\_pow3\_cli\_showcase.sh} compares hybrid and pure
CFF for the six-edge graph with a triple repeated propagator.  The numerator is
\texttt{dot(edges[0], ext[0]) + edges[3][0]}.

\begin{center}
\begin{tabular}{llrrrr}
\toprule
label & backend & orientations & maps & value type & per sample \\
\midrule
hybrid & compiled & 17 & 17 & real & \(0.0955\,\mu s\) \\
hybrid & eager & 17 & 17 & real & \(0.521\,\mu s\) \\
hybrid & eager+symjit & 17 & 17 & real & \(0.104\,\mu s\) \\
hybrid & compiled & 17 & 17 & complex & \(0.271\,\mu s\) \\
hybrid & eager & 17 & 17 & complex & \(0.725\,\mu s\) \\
hybrid & eager+symjit & 17 & 17 & complex & \(0.0989\,\mu s\) \\
CFF & compiled & 62 & 62 & real & \(0.162\,\mu s\) \\
CFF & eager & 62 & 62 & real & \(1.44\,\mu s\) \\
CFF & eager+symjit & 62 & 62 & real & \(0.142\,\mu s\) \\
CFF & compiled & 62 & 62 & complex & \(0.466\,\mu s\) \\
CFF & eager & 62 & 62 & complex & \(2.03\,\mu s\) \\
CFF & eager+symjit & 62 & 62 & complex & \(0.136\,\mu s\) \\
\bottomrule
\end{tabular}
\end{center}

The hybrid representation has fewer numerator maps than CFF, which explains
the faster compiled real evaluator.  Complex compiled evaluation is slower
because the generated function propagates complex arithmetic even though the
input kinematics are real.  Symjit eager evaluation is competitive for this
small expression because the batch is small and the expression has only 17 or
62 maps.

\section{One-Loop 10-Gon}

The script
\texttt{examples/scripts/one\_loop\_10\_external\_runtime\_compare.sh}
compares numerator-free LTD and CFF on a true one-loop 10-gon.

\begin{center}
\begin{tabular}{llrrr}
\toprule
family & backend & orientations & maps & per sample \\
\midrule
LTD & compiled & 10 & 10 & \(0.234\,\mu s\) \\
LTD & eager & 10 & 10 & \(0.852\,\mu s\) \\
LTD & eager+symjit & 10 & 10 & \(0.147\,\mu s\) \\
CFF & compiled & 1022 & 1022 & \(3.96\,\mu s\) \\
CFF & eager & 1022 & 1022 & \(36.6\,\mu s\) \\
CFF & eager+symjit & 1022 & 1022 & \(2.97\,\mu s\) \\
\bottomrule
\end{tabular}
\end{center}

The result is expected: the one-loop CFF orientation count grows like
\(2^n-2\), while LTD has only one cut per internal edge.  The CFF expression is
still fast in absolute terms after compilation, but it is much larger.

\section{Five-Loop Four-External Benchmarks}

The script
\texttt{examples/scripts/five\_loop\_symbolica\_runtime\_compare.sh} compares
three numerator choices on a non-repeated five-loop graph and on the repeated
five-loop graph.  The three numerator choices are \(1\), one edge-external dot
product per internal edge, and the same numerator with three selected edge
terms squared.

\begin{center}
\begin{tabular}{llrrrr}
\toprule
case & family & maps & compiled & eager & eager+symjit \\
\midrule
no repeats, \(N=1\) & LTD & 137 & \(0.326\) & \(1.88\) & \(0.250\) \\
no repeats, \(N=1\) & CFF & 930 & \(1.26\) & \(12.3\) & \(0.892\) \\
no repeats, \(N=1\) & hybrid & 137 & \(0.333\) & \(1.87\) & \(0.246\) \\
no repeats, linear & LTD & 137 & \(0.481\) & \(3.29\) & \(0.379\) \\
no repeats, linear & CFF & 930 & \(1.96\) & \(17.5\) & \(1.35\) \\
no repeats, linear & hybrid & 137 & \(0.484\) & \(3.24\) & \(0.391\) \\
no repeats, squared & LTD & 137 & \(0.493\) & \(3.45\) & \(1.73\) \\
no repeats, squared & CFF & 4692 & \(18.7\) & \(45.9\) & \(13.2\) \\
no repeats, squared & hybrid & 137 & \(0.520\) & \(3.31\) & \(0.390\) \\
\bottomrule
\end{tabular}
\end{center}

\begin{center}
\begin{tabular}{llrrrr}
\toprule
case & family & maps & compiled & eager & eager+symjit \\
\midrule
repeated, \(N=1\) & LTD & 137 & \(0.334\) & \(1.90\) & \(0.272\) \\
repeated, \(N=1\) & CFF & 930 & \(1.26\) & \(12.5\) & \(0.899\) \\
repeated, \(N=1\) & hybrid & 630 & \(0.823\) & \(4.88\) & \(0.579\) \\
repeated, linear & LTD & 137 & \(0.481\) & \(3.44\) & \(0.408\) \\
repeated, linear & CFF & 930 & \(1.88\) & \(17.4\) & \(1.32\) \\
repeated, linear & hybrid & 630 & \(21.4\) & \(37.8\) & \(12.4\) \\
repeated, squared & LTD & 137 & \(0.491\) & \(3.53\) & \(0.387\) \\
repeated, squared & CFF & 115 & \(1.89\) & \(9.49\) & \(1.16\) \\
repeated, squared & hybrid & 115 & \(1.84\) & \(9.47\) & \(1.11\) \\
\bottomrule
\end{tabular}
\end{center}

For non-repeated graphs the hybrid family collapses to the LTD structure, so
their orientation and map counts agree.  CFF is slower because it has more
causal sign sectors.  The squared CFF no-repeat case is much larger because the
bounded-degree contact completion adds pinched CFF minors.

For repeated graphs, the LTD rows are included only as runtime references.  The
ordinary unsplit LTD expression does not include numerator derivatives for
raised propagators; the correct comparison is CFF versus hybrid versus the
split-mass LTD limiting test.  The repeated unit numerator shows the expected
middle ground: hybrid has fewer maps than CFF but more than LTD.  The repeated
linear case is slower for hybrid because the confluent interpolation creates
heavier denominator variants per numerator map.  In the squared repeated case,
the bounded-degree reduction compresses both CFF and hybrid to 115 maps and the
two timings become nearly identical.

\begin{thebibliography}{9}
  \bibitem{ltd1906}
  J. J. Aguilera-Verdugo et al.,
  algorithmic multiloop LTD construction, arXiv:1906.06138.

  \bibitem{cff2211}
  Z. Capatti et al.,
  causal/CFF representation construction, arXiv:2211.09653.

  \bibitem{cff2311}
  Z. Capatti et al.,
  further CFF representation and box examples, arXiv:2311.14374.

  \bibitem{cltd}
  \texttt{cLTD} and \texttt{gammaLoop} implementations,
  \url{https://github.com/alphal00p/cLTD} and
  \url{https://github.com/alphal00p/gammaloop}.
\end{thebibliography}

\end{document}
